\documentclass[10pt]{article}
\usepackage[margin=0.78in]{geometry}
\usepackage{amsmath,amssymb,amsthm}
\usepackage{microtype}
\usepackage[hidelinks]{hyperref}
\newtheorem*{answer}{Answering theorem}
\begin{document}
\begin{center}
{\Large\bfseries Literature resolution: self-adjoint finite-dimensional\par
Lindbladians have positive complete log-Sobolev constant\par}
\medskip
Automated Math Discovery Run \texttt{fa\_banach\_001}; Model: GPT5.6\par
9 August 2026
\end{center}
\medskip

\begin{center}
\fbox{\begin{minipage}{0.9\textwidth}
\textbf{Status: literature already answered (Conjecture 1.4 only).}
The conjecture in arXiv:2006.14578 that every self-adjoint Lindbladian
$L$ satisfies $\operatorname{CLSI}(L)>0$ follows from Theorem~3.3 of
Gao--Rouz\'e, arXiv:2102.04146.  The published source paper explicitly
records this later resolution.
\end{minipage}}
\end{center}

\section*{Source conjecture}
Junge, Li, and LaRacuente formulate the following as Conjecture~1.4:
\[
  \text{for every self-adjoint Lindbladian }L,
  \qquad \operatorname{CLSI}(L)>0.
\]
It appears on printed p.~4 of the arXiv source paper
\cite[p.~4, Conjecture~1.4]{JLL}.  Here ``self-adjoint'' is with respect to
the tracial/GNS inner product, and the ambient algebra in this Lindbladian
setting is a finite-dimensional matrix algebra.

\section*{Exact later theorem}
\begin{answer}[Gao--Rouz\'e, Theorem 3.3]
Let $(P_t=e^{tL})_{t\geq0}$ be a GNS-symmetric quantum Markov semigroup on a
finite-dimensional von Neumann algebra $\mathcal M\subseteq B(H)$, and suppose
it extends to a semigroup on $B(H)$.  If $\mathcal F$ is its fixed-point
algebra, then it satisfies the complete modified logarithmic Sobolev
inequality and
\[
 \frac{\lambda(L)}{C_{\tau,\mathrm{cb}}(\mathcal M:\mathcal F)}
 \leq \alpha_{\mathrm{CMLSI}}(L)\leq 2\lambda(L).
\]
\end{answer}
This is Theorem~3.3 on printed p.~20 of arXiv:2102.04146
\cite[p.~20, Theorem~3.3]{GR}.  In the source setting
$\mathcal M=B(H)=M_n$, so the extension hypothesis is automatic.
Self-adjointness for the trace is precisely GNS symmetry for the maximally
mixed invariant state.  Thus the theorem applies to every source
Lindbladian and yields the claimed complete inequality; in the usual
nontrivial quotient by the fixed algebra the spectral gap is positive by
finite-dimensionality.  Degenerate fixed directions are already removed by
the conditional expectation onto $\mathcal F$.

The identification is not merely terminological.  The 2025 published
version of \emph{Graph H\"ormander Systems} explicitly states that Gao and
Rouz\'e subsequently proved the conjectured positivity.

\section*{Scope boundary: a different conjecture remains open}
The same source also asks for an anti-transference identity recovering
$\operatorname{CLSI}(\Delta_X)$ as an infimum of constants of
finite-dimensional transferred generators.  The published paper proves the
corresponding identity for the strengthened $\operatorname{CLSI}^{+}$
constant but continues to label the plain-$\operatorname{CLSI}$ identity a
conjecture.  Gao--Rouz\'e's finite-dimensional positivity theorem does not
establish that equality.  Accordingly, this packet closes only
Conjecture~1.4 and makes no claim that every open problem in arXiv:2006.14578
has been solved.

\begin{thebibliography}{9}
\bibitem{JLL}
M. Junge, H. Li, and N. LaRacuente,
\emph{Graph H\"ormander Systems},
Ann. Henri Poincar\'e 26 (2025), 2683--2736;
\href{https://arxiv.org/abs/2006.14578}{arXiv:2006.14578}.

\bibitem{GR}
L. Gao and C. Rouz\'e,
\emph{Complete entropic inequalities for quantum Markov chains},
Arch. Ration. Mech. Anal. 245 (2022), 183--238;
\href{https://arxiv.org/abs/2102.04146}{arXiv:2102.04146}.
\end{thebibliography}
\end{document}
